% p1-04-target-dataset

\section{P1 §4. Target Dataset and Pre-Registration Freeze Protocol}

4. Target Dataset and Pre-Registration Freeze Protocol

                   Figure (fig-p1-ch4-validation). Validation triptych --- Target / Experimental Band /
                                                 Pre-Registration Lock.

  4.1 Dataset Selection Criteria
  The target dataset T consists of dimensionless physical constants drawn from standard,
  publicly available experimental compilations. A dimensionless constant is defined as a
  ratio of physical quantities carrying the same dimension, so that its numerical value is
  independent of any choice of units. Only dimensionless quantities are included;
  dimensional constants (e.g., Newton's gravitational constant G in SI units) depend on
  an arbitrary unit choice and are excluded.

  Candidate constants for inclusion must satisfy all of the following criteria:
  1. Dimensionless (ratio of like-dimensioned quantities).
  2. Drawn from the Standard Model or cosmological concordance model.
  3. Tabulated in the PDG Review of Particle Physics 2024 \citep{pdg2024} or the CODATA 2022 \citep{codata2022}
  recommended values.
  4. Relative experimental uncertainty better than 10-2.

  5. Not selected on the basis of proximity to any symbolic expression.

  4.2 Candidate Constants
  The following candidate constants meet the selection criteria. All values are taken from
  the PDG 2024 Review of Particle Physics \citep{pdg2024} and the CODATA 2022 recommended values \citep{codata2022};
  cosmological parameters follow the Planck 2018 cosmological parameter results \citep{planck2018vi} as
  adopted by the PDG.

  Constants with relative uncertainty > 10-2 are retained as secondary targets but are
  not included in the primary compressibility test; they appear in Section 10 as part of
  the negative control analysis.

  4.3 Pre-Registration Freeze Protocol
  To prevent post-hoc selection bias --- specifically, the practice of including in T only
  those constants for which the grammar achieves small residuals --- the dataset must be
  frozen prior to any symbolic search. The following protocol is adopted and constitutes a
  formal pre-registration commitment.

  Protocol P1 (Dataset Freeze):

  1. The set T is defined by listing all dimensionless parameters satisfying the criteria in
  Section 4.1, from the PDG Review of Particle Physics 2024 \citep{pdg2024} and CODATA 2022 \citep{codata2022}, with
  PDG edition fixed at the 2024 release (Phys. Rev. D 110, 030001).
  2. This list is committed to a version-controlled repository with a cryptographic hash,
  before any search computation is performed.
  3. No constant may be added to T after any symbolic search has been performed,
  except as part of a prospectively defined out-of-sample test (Section 10).
  4. If a constant's experimental value is updated after the dataset freeze, the updated
  value is used only in the out-of-sample falsification analysis (Section 10.2), not in the
  primary fit analysis.
  5. All data files, version hashes, and Zenodo provenance records are deposited as
  documented in the PhD monograph's Data Availability appendix (PhD: App.G Data
  Availability) and the CLARA Provenance Audit 2026-05-12.

  Rationale. This protocol maps onto standard pre-registration norms in empirical
  science. Post-hoc data selection is among the most potent sources of false discovery in
  symbolic search studies; any violation of Protocol P1 invalidates the statistical
  interpretation of reported residuals.

  4.4 Logarithmic Distribution of the Dataset
  In log-space, the target constants {log Xi} span approximately four decades (from
  log(0.022) to log(137)) and are irregularly distributed. No claim of log-uniformity is
  made. The empirical log-space distribution is an explicit input to the null-model
  construction (Section 5.2), and any compressibility claim must be evaluated relative to
  this empirical distribution.

